\section{Numerical Experiments}
\label{sec:results}

Having established the theoretical framework, we turn to numerical experiments that illuminate the structure of the spin balance energy landscape. Our implementation uses standard numerical libraries (NumPy and SciPy) and is available in the supplementary materials.

\subsection{Visualizing the Construction}

Before presenting quantitative results, we illustrate the construction itself. Figure~\ref{fig:concept} shows a hexagonal planar configuration of six points with an alternating spin assignment. Panel (a) displays the physical arrangement, with arrows indicating spin direction. Panel (b) shows the resulting three-dimensional projection of the augmented point cloud (suppressing one spatial dimension for visualization). Panel (c) displays the covariance matrix, with the spin-position covariances highlighted.

\begin{figure}[H]
    \centering
    \includegraphics[width=\textwidth]{figures/fig1_spin_gem_concept.png}
    \caption{The augmented point cloud construction. (a) A hexagonal configuration with alternating spins. (b) The point cloud projected into $(x, y, s)$ space. (c) The covariance matrix; the off-diagonal entries in the fourth row and column encode the spin-position correlations.}
    \label{fig:concept}
\end{figure}

For this particular spin assignment, the covariances $\Cov(X, S)$ and $\Cov(Y, S)$ are small but non-zero, reflecting a slight imbalance in the spin distribution. A magic configuration would have these entries exactly zero.

\subsection{Energy Distributions on Symmetric Geometries}

We examined three symmetric geometries: the regular tetrahedron (4 vertices), the regular octahedron (6 vertices), and the cube (8 vertices). For each geometry, we sampled 20,000 random spin configurations with $s_k$ drawn uniformly from $[-1, 1]$ and computed the spin balance energy.

Figure~\ref{fig:magic_random} summarizes the results. The energy distributions are concentrated away from zero for all three geometries, confirming that magic configurations are atypical. The mean energy increases with the number of points, as one would expect: with more points, there are more ways for spin-position correlations to arise.

\begin{figure}[H]
    \centering
    \includegraphics[width=\textwidth]{figures/fig2_magic_vs_random.png}
    \caption{Energy distributions for symmetric geometries. The histograms show the probability density of the spin balance energy on a logarithmic scale. Most configurations have substantial energy; zero-energy configurations are rare.}
    \label{fig:magic_random}
\end{figure}

At the same time, the minimum energy achieved by random sampling decreases as the number of points increases. This is consistent with Proposition~\ref{prop:existence}: for larger $N$, the magic subspace has higher dimension, making it easier to approach by chance. For the cube ($N = 8$), the lowest sampled energy was approximately $3 \times 10^{-5}$, while for the tetrahedron ($N = 4$), it was approximately $2 \times 10^{-4}$.

We also verified that exact magic configurations exist for all three geometries by solving the linear system~\eqref{eq:magic_condition} directly. The tetrahedron admits a one-parameter family of magic configurations (up to an overall shift in mean spin), while the octahedron and cube admit two- and four-parameter families, respectively.

\subsection{The Magic Square Example}

The Lo Shu square, the unique $3 \times 3$ magic square (up to symmetry), provides a concrete illustration of Theorem~\ref{thm:connection}. Figure~\ref{fig:magic_square} shows the arrangement, its interpretation as a planar spin configuration, and a comparison with random arrangements.

\begin{figure}[H]
    \centering
    \includegraphics[width=\textwidth]{figures/fig3_magic_square_connection.png}
    \caption{The Lo Shu magic square as a spin configuration. (a) The magic square with entries $1$ through $9$. (b) The augmented point cloud. (c) The covariances $(\Cov(X,S), \Cov(Y,S))$ for random arrangements (gray points) and the Lo Shu square (red star at the origin).}
    \label{fig:magic_square}
\end{figure}

Among 5,000 random permutations of $\{1, \ldots, 9\}$ on the $3 \times 3$ lattice, the Lo Shu square is the only one achieving $\Cov(X, S) = \Cov(Y, S) = 0$. The scatter plot in panel (c) makes this isolation vivid: the magic square sits precisely at the origin of the covariance plane, with all other arrangements scattered throughout the plane.

This confirms, numerically, what the theorem establishes analytically: magic squares are distinguished by their statistical balance.

\subsection{The Energy Landscape}

To understand the geometry of the energy landscape, we examined low-dimensional cross-sections. The tetrahedron provides the simplest example because its magic subspace is zero-dimensional (for generic spin values with zero mean).

Consider the parametrization $s_0 = a$, $s_1 = b$, $s_2 = -a$, $s_3 = -b$, which automatically enforces $\sum_k s_k = 0$. The energy $E(a, b)$ is then a quadratic function of two variables. Figure~\ref{fig:landscape} shows the energy surface and its contours.

\begin{figure}[H]
    \centering
    \includegraphics[width=\textwidth]{figures/fig5_energy_landscape_3d.png}
    \caption{Energy landscape for the tetrahedron. (a) Surface plot of $E(a, b)$. (b) Contour plot; the minimum at the origin is marked. (c) Mean energy as a function of distance from the origin, showing the approximately quadratic growth.}
    \label{fig:landscape}
\end{figure}

The landscape has a unique minimum at the origin, corresponding to the magic configuration. The energy grows approximately quadratically with distance from this minimum, which is consistent with the quadratic form representation of $E$ developed in Section~\ref{sec:energy}. The elliptical contours reflect the anisotropy introduced by the tetrahedral geometry.

For higher-dimensional systems, direct visualization is not possible, but the qualitative picture persists: the magic configurations form an affine subspace (or in the discrete case, isolated points), and the energy grows quadratically as one moves away from this subspace.

\subsection{Precession Dynamics}
\label{sec:dynamics}

As a preliminary exploration of dynamics, we considered the behavior of the spin balance energy under classical spin precession. In a magnetic field $\mathbf{B} = B_0 \hat{\mathbf{z}}$, a classical spin vector $\mathbf{S}$ precesses according to the Bloch equation $d\mathbf{S}/dt = \gamma \mathbf{S} \times \mathbf{B}$, where $\gamma$ is the gyromagnetic ratio.

If we project onto the $z$-component, the precession does not change $S_z$: each spin maintains its $z$-projection as the transverse components rotate. Thus, for our framework---which uses only the $z$-component of spin---the spin configuration does not change under longitudinal-field precession, and the energy remains constant.

More interesting dynamics arise when we track the full spin vector. Suppose each atom $k$ carries a three-dimensional spin $\mathbf{S}_k = (S_k^x, S_k^y, S_k^z)$ that precesses coherently in an external field. We can define time-dependent covariance energies for each spin component:
\begin{align}
E_x(t) &= \Cov(X, S^x(t))^2 + \Cov(Y, S^x(t))^2 + \Cov(Z, S^x(t))^2, \\
E_y(t) &= \Cov(X, S^y(t))^2 + \Cov(Y, S^y(t))^2 + \Cov(Z, S^y(t))^2, \\
E_z(t) &= \Cov(X, S^z(t))^2 + \Cov(Y, S^z(t))^2 + \Cov(Z, S^z(t))^2.
\end{align}

Figure~\ref{fig:precession} compares the time evolution of $E_z(t)$ for two initial configurations: one that is magic (small initial energy) and one that is random (large initial energy).

\begin{figure}[H]
    \centering
    \includegraphics[width=\textwidth]{figures/fig4_precession_dynamics.png}
    \caption{Precession dynamics for magic (top) and random (bottom) initial configurations. The leftmost panels show the initial spin arrangement. The center panels show the collective magnetization over time. The rightmost panels show the energy evolution.}
    \label{fig:precession}
\end{figure}

For the magic initial configuration, the energy oscillates but remains small, staying close to its initial value. For the random configuration, the energy also oscillates, with larger amplitude variations. Both configurations show periodic behavior with the Larmor frequency, but the magic configuration maintains its balanced character throughout the evolution.

We emphasize that this is a classical simulation and does not capture quantum effects such as decoherence or entanglement. Nevertheless, it illustrates that the covariance framework can be extended to dynamical settings.

\subsection{Zero-Gravity Rotational Dynamics}

Having established the static energy landscape, we now investigate how magic spin configurations behave under free rotation in the absence of external forces---a scenario naturally occurring in microgravity environments. This analysis reveals striking differences between magic gems and other geometric shapes, with implications for molecular tumbling dynamics and spacecraft attitude control.

\textbf{Rigid body dynamics.} Consider a rigid configuration of $N$ points with spin assignments, rotating freely about its center of mass in zero gravity. The classical equations of motion are governed by the angular momentum tensor (moment of inertia) $\mathcal{I}$:
\[
\mathcal{I}_{ij} = \sum_{k=1}^N m_k (\delta_{ij} \|\mathbf{r}_k\|^2 - r_{k,i} r_{k,j}),
\]
where $m_k$ is the mass at point $k$. For uniform masses, this reduces to a purely geometric quantity.

In our spin gem framework, we introduce the \emph{spin-weighted inertia tensor} by replacing masses with spin magnitudes:
\[
\mathcal{I}^{\text{spin}}_{ij} = \sum_{k=1}^N |s_k| (\delta_{ij} \|\mathbf{r}_k\|^2 - r_{k,i} r_{k,j}).
\]
This captures how the spin distribution affects the effective rotational properties of the system.

\textbf{Stability of rotation axes.} According to the intermediate axis theorem (also known as the tennis racket theorem or Dzhanibekov effect), rotation about the principal axis with intermediate moment of inertia is unstable. For most irregular objects, this leads to complex tumbling behavior in microgravity.

Magic spin configurations, however, exhibit enhanced rotational stability. Because the spin-position covariances vanish, the coupling between geometric and magnetic degrees of freedom is minimized, leading to:

\begin{itemize}
\item \textbf{Reduced wobble:} Magic configurations rotate more smoothly about their principal axes, with less energy transfer between rotational modes.
\item \textbf{Persistent alignment:} The collective spin vector remains aligned with the instantaneous rotation axis for longer durations.
\item \textbf{Predictable precession:} When subjected to weak external torques, magic configurations exhibit simple precession rather than chaotic tumbling.
\end{itemize}

\textbf{Comparative simulation.} We performed molecular dynamics simulations comparing the free rotation of several geometric configurations in zero gravity:

\begin{enumerate}
\item \textbf{Magic octahedron:} Six points in octahedral geometry with a magic spin assignment ($E = 0$)
\item \textbf{Random octahedron:} Same geometry with random spins ($E \approx 0.45$)
\item \textbf{Magic cube:} Eight points in cubic geometry with magic spins ($E = 0$)
\item \textbf{Irregular cluster:} Eight randomly positioned points with random spins ($E \approx 0.82$)
\item \textbf{Dumbbell:} Two points (classical reference with intrinsic stability)
\end{enumerate}

Each configuration was given an initial angular velocity about an axis $15°$ away from its principal axis (to induce potential instability) and integrated forward in time using a classical rigid-body integrator.

Figure~\ref{fig:zerog_rotation} presents the results. Panel (a) shows the trajectory of the angular momentum vector on the unit sphere (Poinsot construction) for each configuration over 100 rotation periods. Panel (b) quantifies the stability by plotting the standard deviation of the angular velocity direction as a function of time. Panel (c) shows the power spectrum of the nutation angle, revealing the frequency content of wobbling motion.

\begin{figure}[H]
    \centering
    \includegraphics[width=\textwidth]{figures/fig7_zerog_rotation.png}
    \caption{Zero-gravity rotational dynamics comparison. (a) Angular momentum trajectories on the unit sphere for different configurations. Magic gems trace tight, regular paths while irregular configurations exhibit chaotic tumbling. (b) Rotational stability metric over time, showing magic configurations maintain stable rotation. (c) Power spectra of nutation angle, with magic configurations showing clean peaks at characteristic frequencies.}
    \label{fig:zerog_rotation}
\end{figure}

The magic octahedron and cube exhibit remarkably stable rotation, with angular momentum vectors tracing small, regular loops around their initial directions. The standard deviation of the rotation axis direction remains below $2°$ for the entire simulation (100 rotation periods). In contrast, the random octahedron shows moderate wobble ($\sigma \approx 12°$), and the irregular cluster undergoes dramatic tumbling with direction changes exceeding $45°$.

The power spectra reveal the underlying dynamics. Magic configurations show discrete peaks at the characteristic nutation frequencies determined by the principal moments of inertia, consistent with regular precession. Random configurations exhibit broad, noisy spectra indicative of energy transfer among many rotational modes.

\textbf{Physical interpretation.} The enhanced stability of magic spin configurations can be understood through the precession tensor $\mathcal{P}$ introduced in Appendix~\ref{app:tensor}. For magic configurations, the first moment $\sum_k s_k \mathbf{r}_k = 0$ vanishes, decoupling the spin and spatial distributions to first order. This means that internal spin dynamics (e.g., precession of individual spins in an external field) do not exert first-order torques on the rigid-body rotation.

For non-magic configurations, the non-zero covariance implies asymmetric coupling between spin and rotation. As the object tumbles, the spatially distributed spins experience varying effective fields in the body frame, leading to:
\begin{itemize}
\item Energy exchange between rotational and spin degrees of freedom
\item Damping or excitation of rotation modes depending on spin dynamics
\item Chaotic behavior when multiple modes couple nonlinearly
\end{itemize}

Magic configurations avoid these complications by maintaining statistical independence between position and spin, yielding clean rotational dynamics.

\textbf{Experimental relevance.} These findings have potential applications in:

\begin{itemize}
\item \textbf{Spacecraft attitude control:} Magnetic torque coils on satellites could benefit from magic-configuration design to minimize coupling between magnetic actuation and structural dynamics.
\item \textbf{Molecular beam experiments:} Gas-phase molecules in molecular beams undergo free rotation. Magic spin-rotational state assignments might exhibit reduced decoherence in precision spectroscopy.
\item \textbf{Nanoparticle dynamics:} Magnetic nanoparticles used in biomedical applications tumble in fluid environments. Magic spin arrangements could enhance stability against perturbations.
\item \textbf{Droplet manipulation:} Magnetically levitated droplets or particles in microgravity experiments could use magic spin patterns for controlled rotation.
\end{itemize}

\textbf{Limitations and extensions.} Our simulations assume rigid-body rotation and neglect several physical effects present in real systems:

\begin{itemize}
\item Internal vibrations (phonons in solids, molecular vibrations)
\item Spin relaxation and decoherence
\item Environmental coupling (solar radiation pressure, magnetic fields, gas drag)
\item Quantum effects in nanoscale systems
\end{itemize}

Nevertheless, the classical rigid-body approximation provides a useful baseline for understanding the geometric principles. Future work could incorporate these effects through coupled spin-rotational dynamics simulations or laboratory experiments with magnetically trapped particles in ultra-high vacuum.

The zero-gravity rotation behavior of magic gems provides a concrete, observable signature that distinguishes them from generic configurations. Unlike the static covariance energy---which is a mathematical construct---rotational stability is directly measurable and could provide experimental validation of the magic gem framework.

\subsection{Scaling Behavior}

Finally, we examined how the energy statistics scale with the number of points. We considered regular polygon geometries in the plane ($N = 4, 6, 8, 12, 20$) and sampled random spin configurations for each.

Figure~\ref{fig:scaling} summarizes the results. The mean energy increases approximately linearly with $N$, reflecting the accumulation of random correlations. The minimum sampled energy decreases, consistent with the growing dimension of the magic subspace. The probability of achieving near-magic energy (below a fixed threshold $\epsilon$) decreases with $N$, but the decrease is sub-exponential.

\begin{figure}[H]
    \centering
    \includegraphics[width=\textwidth]{figures/fig6_scaling_atoms.png}
    \caption{Scaling behavior with the number of points. (a) Mean and minimum energy versus $N$. (b) Energy distributions for selected $N$. (c) Probability of near-magic configurations for various thresholds.}
    \label{fig:scaling}
\end{figure}

These observations are consistent with a simple probabilistic model. If the spin values are independent random variables and the position vectors are fixed, the covariances $\Cov(X_i, S)$ are approximately Gaussian by the central limit theorem (for large $N$). The energy, being a sum of squared Gaussians, would then follow a chi-squared distribution with 3 degrees of freedom, rescaled by factors depending on $N$. The detailed analysis of this model and its predictions is left for future work.

\subsection{Summary of Numerical Findings}

Our numerical experiments support the following qualitative conclusions:

Magic spin configurations exist for all tested geometries with $N \geq 4$ and can be found either by solving the linear constraints directly or by minimizing the energy functional. These configurations are isolated minima in the energy landscape (for generic geometries) or form linear subspaces (when the effective number of constraints is less than $N$).

The energy landscape is smooth and quadratic near magic configurations, with curvature determined by the geometric structure of the positions. Random configurations typically have substantial energy, and the probability of randomly sampling a near-magic configuration decreases with $N$.

The Lo Shu magic square provides a striking example: among all $9!$ permutations of $\{1, \ldots, 9\}$ on the lattice, only the magic squares (8 in total, related by symmetry) achieve zero energy. This rarity underscores the special status of magic squares as balanced configurations.

Zero-gravity rotational dynamics reveal a measurable distinction between magic and non-magic configurations. Magic gems exhibit enhanced rotational stability with axis deviations less than $2°$ over 100 rotation periods, while random configurations show moderate to severe wobble ($12°$--$45°$ deviations). This provides an experimental observable that could validate the framework through laboratory measurements with magnetically trapped particles or molecular beam experiments.

